\subsectionaddtolist{الف}

{جهت اول (از راست به چپ):} فرض کنید یک ماشین شمارشگر مانند $E$ داریم که رشته‌های یک زبان مانند $L$ را به ترتیب متعارف چاپ می‌کند. می‌خواهیم نشان دهیم $L$ تصمیم‌پذیر است.
برای این کار یک ماشین تصمیم‌گیر $M$ برای $L$ به این صورت می‌سازیم:

ماشین $M$ روی ورودی $w$:
\begin{enumerate}
    \item ماشین شمارشگر $E$ را شبیه‌سازی می‌کند.
    \item برای هر رشته $s$ که توسط $E$ چاپ می‌شود، آن را با $w$ مقایسه می‌کند.
    \item اگر $s = w$ بود، ماشین $M$ ورودی را می‌پذیرد.
    \item اگر $s$ در ترتیب متعارف از $w$ بزرگتر بود، با توجه به اینکه $E$ رشته‌ها را به صورت صعودی چاپ می‌کند، ماشین $M$ مطمئن می‌شود که $w$ هرگز چاپ نخواهد شد. پس شبیه‌سازی را متوقف کرده و ورودی را رد می‌کند.
    \item اگر ماشین $E$ متوقف شد و تا آن لحظه $w$ چاپ نشده بود، $M$ ورودی را رد می‌کند (این حالت برای زبان‌های متناهی رخ می‌دهد).
\end{enumerate}
چون طول رشته ورودی متناهی است، $M$ همواره در زمان متناهی (یا با پیدا کردن $w$، یا با رسیدن به رشته‌ای بزرگتر از آن، یا با پایان یافتن کار $E$) متوقف می‌شود، پس $L$ تصمیم‌پذیر است.

{جهت دوم (از چپ به راست):} فرض کنید زبان $L$ تصمیم‌پذیر است و ماشین تصمیم‌گیر $M$ آن را می‌پذیرد. می‌خواهیم یک ماشین شمارشگر $E$ بسازیم که رشته‌های آن را به ترتیب متعارف چاپ کند.

نحوه کار ماشین $E$:
\begin{enumerate}
    \item رشته‌های $\Sigma^*$ را به ترتیب متعارف تولید می‌کند (مانند $\epsilon, 0, 1, 00, 01, \dots$).
    \item هر رشته تولید شده مانند $s$ را به عنوان ورودی به ماشین $M$ می‌دهد.
    \item چون $M$ تصمیم‌گیر است، روی هر ورودی متوقف می‌شود. اگر $M$ رشته $s$ را پذیرفت، ماشین $E$ رشته $s$ را چاپ می‌کند و در غیر این صورت از آن عبور می‌کند.
\end{enumerate}
به این ترتیب، ماشین $E$ تمام رشته‌های ممکن را به ترتیب بررسی کرده و فقط رشته‌های متعلق به زبان $L$ را دقیقاً به همان ترتیب متعارف چاپ می‌کند.

\vspace{0.5cm}

\subsectionaddtolist{ب}

فرض کنید $L$ یک زبان تشخیص‌پذیر و نامتناهی باشد. از آنجا که $L$ تشخیص‌پذیر است، پس یک ماشین شمارشگر مانند $E_L$ دارد که رشته‌های آن را چاپ می‌کند (این ماشین لزوماً رشته‌ها را به ترتیب متعارف چاپ نمی‌کند).
ما یک ماشین شمارشگر جدید مانند $E'$ می‌سازیم که زیرمجموعه‌ای از رشته‌های چاپ شده توسط $E_L$ را انتخاب کرده و چاپ می‌کند.

نحوه کار ماشین $E'$:
\begin{enumerate}
    \item ماشین $E_L$ را شبیه‌سازی می‌کند.
    \item اولین رشته‌ای که توسط $E_L$ چاپ می‌شود را به عنوان خروجی خود چاپ کرده و آن را در متغیری به نام $w_{max}$ ذخیره می‌کند.
    \item شبیه‌سازی $E_L$ را ادامه می‌دهد. هر بار که $E_L$ رشته جدیدی مانند $w$ چاپ کرد، ماشین $E'$ بررسی می‌کند که آیا $w$ در ترتیب متعارف اکیداً بزرگتر از $w_{max}$ است یا خیر.
    \item اگر $w > w_{max}$ بود، ماشین $E'$ رشته $w$ را چاپ می‌کند و مقدار $w_{max}$ را برابر با $w$ قرار می‌دهد. در غیر این صورت، از آن رشته صرف نظر می‌کند.
\end{enumerate}

با این روش، ماشین $E'$ رشته‌ها را همواره به صورت صعودی و با ترتیب متعارف چاپ می‌کند.
از آنجا که زبان $L$ نامتناهی است، ماشین $E_L$ در نهایت بی‌نهایت رشته متمایز را چاپ خواهد کرد. بنابراین، ماشین $E'$ در طول شبیه‌سازی خود همواره در مقطعی به رشته‌ای برمی‌خورد که از $w_{max}$ فعلی بزرگتر است. در نتیجه روند چاپ متوقف نخواهد شد و ماشین $E'$ بی‌نهایت رشته را چاپ می‌کند.

زبان تولید شده توسط ماشین $E'$ را $A$ می‌نامیم. چون تمام رشته‌های $A$ توسط $E_L$ نیز تولید شده‌اند، واضح است که $A \subseteq L$ برقرار است. همچنین از آنجا که $A$ نامتناهی است و ماشین $E'$ رشته‌های آن را دقیقاً به ترتیب متعارف چاپ می‌کند، طبق قضیه ثابت شده در قسمت الف، زبان $A$ یک زبان تصمیم‌پذیر است.
در نتیجه $A$ یک زیرمجموعه نامتناهی و تصمیم‌پذیر برای $L$ است.
